\section{Technologies used}
\label{sec:tech}

\subsection{The stack, and why each piece is there}

\begin{xltabular}{\linewidth}{L{4.3cm} L{2.15cm} X}
\toprule
\textbf{Technology} & \textbf{Version} & \textbf{Why this one} \\
\midrule
\endfirsthead
\toprule
\textbf{Technology} & \textbf{Version} & \textbf{Why this one} \\
\midrule
\endhead
Kotlin & 1.9.24 &
Primary language: coroutines and \code{Flow} give the whole app one asynchronous idiom, and Room's
suspend/Flow DAOs mean no database call ever needs a manual thread hop. \\
\addlinespace
Java & 17 target &
Used deliberately for framework-free logic in \code{util/}: parsers, the force-directed layout, the
streaming envelope reader, backup writer/reader and the diff. Pure Java classes with no Android imports
are the parts worth unit-testing on the JVM, so the language boundary and the test boundary coincide. \\
\addlinespace
Android SDK (XML Views) & min 24,\newline target 35 &
Views rather than Compose: the design reference is a Views project, minSdk 24 keeps the reach wide, and
the two custom-drawing requirements (the graph canvas, the reading column) are more direct against
\code{View} and \code{Canvas}. \\
\addlinespace
AGP / Gradle & 8.7.3 / 8.9 &
Current stable pair for compileSdk 35; the wrapper is committed so the build needs no global Gradle. \\
\addlinespace
Room & 2.6.1 &
Local source of truth. Chosen over raw SQLite or files for four properties this project depends on:
compile-time-checked queries, real transactions, observable \code{Flow} queries, and testable migrations.
FTS4 comes from Room's \code{@Fts4} annotation --- version 2.6 has no FTS5 annotation, which is why the
index is FTS4. \\
\addlinespace
KSP & 1.9.24-1.0.20 &
Room's annotation processor, roughly twice as fast as kapt on this project. \\
\addlinespace
Navigation Component & 2.8.3 &
One activity, five bottom-navigation tabs and nine pushed destinations with typed arguments and a real
back stack, all declared in \code{nav\_graph.xml}. Global actions keep cross-screen navigation from
turning into a web of per-screen actions. \\
\addlinespace
Lifecycle / ViewModel & 2.8.7 &
\code{ViewModel} + \code{StateFlow} + \code{repeatOnLifecycle} is the app's entire state-management
story; every screen exposes one immutable UI-state data class. \\
\addlinespace
Material Components & 1.12.0 &
Chips, bottom sheets, dialogs and the bottom navigation bar, all restyled onto the mocha palette rather
than used with default Material colours. \\
\addlinespace
Markwon & 4.6.2 (core, linkify) &
Renders Markdown into a \code{TextView} entirely offline. Its pluggable link resolver is what makes
\code{[[wiki-links]]} clickable: they are rewritten to an internal \code{mocha://post/<id>} scheme before
rendering and intercepted on tap. Only the artefacts actually used are pulled in. \\
\addlinespace
Retrofit + OkHttp + Gson & 2.11.0 / 4.12.0 / 2.11.0 &
Retrofit for the two request/response endpoints, raw OkHttp for the Server-Sent-Events stream (Retrofit
has no streaming idiom worth the ceremony here), Gson for the action envelope. All three are used only
for the gateway --- there is no other network traffic in the app. \\
\addlinespace
kotlinx-coroutines & 1.9.0 &
\code{Dispatchers.IO} for database and network work, \code{Dispatchers.Default} for graph layout, and
structured cancellation so leaving a screen cancels an in-flight AI stream. \\
\addlinespace
Node.js + Express & 24 / 4 &
The gateway. Roughly 390 lines across \code{server.js} and \code{prompts.js}, no database, no state: it
holds the API key, builds the system prompt for the requested mode, normalises upstream errors, and
relays the token stream. \\
\addlinespace
DeepSeek API & \code{deepseek-chat} &
The language model, reached only through the gateway. The app never talks to it directly and never holds
a key. The action protocol is plain JSON documented in the prompt, so swapping providers means changing
one URL. \\
\addlinespace
JUnit 4 + Room testing & 4.13.2 / 2.6.1 &
JVM unit tests for pure logic, plus \code{MigrationTestHelper} for real schema migrations with real rows. \\
\bottomrule
\end{xltabular}

\subsection{Deliberate non-dependencies}

What was left out matters as much as what went in, because every dependency is a thing that can break the
build a day before a deadline:

\begin{itemize}
  \item \textbf{No dependency-injection framework.} Hilt or Koin would earn their keep in a larger team
        project; here the object graph is eight lazily-created singletons on the \code{Application}, which
        is smaller and needs no annotation processor.
  \item \textbf{No graph library.} The knowledge graph is a \code{View} plus about 150 lines of
        Fruchterman--Reingold in Java. A charting or graph SDK would have added megabytes and a WebView
        for less control.
  \item \textbf{No YouTube player SDK.} Playback uses the standard IFrame embed in a locked-down
        \code{WebView} against \code{youtube-nocookie.com}, with the URL allow-list enforced in the
        \code{WebViewClient}.
  \item \textbf{No image loading, no analytics, no crash reporting, no ads.} Nothing in the app needs to
        phone home, and adding an SDK that could would contradict the privacy promise made on the
        settings screen.
  \item \textbf{No Compose.} Mixing Compose into a Views codebase mid-project buys nothing the deadline
        can spend.
\end{itemize}

\subsection{Development tooling}

The repository carries an \code{AGENTS.md} that states the architecture rules a contributor (human or
AI coding assistant) must follow: the layer boundaries, the Kotlin/Java split, the design-token rule, the
phase order, and the invariant that every AI-originated \emph{write} goes through \code{ActionExecutor}
--- read-only context lookups in \code{ContextTools} do query the DAOs directly, but nothing in the AI
subsystem may write to one. \code{docs/plan.md} is the
long-form design document, updated at the end of every phase with notes for the next one. Those files are
development tooling, not part of the shipped application.
